\section{归纳偏置何时帮忙何时挡路}
\label{sec:14-Deep-Learning/04-Convolutional-Networks/05}

你手上有 8000 张标注好的工业缺陷图。同事发来一篇论文，图表显示 Vision Transformer 在大规模基准上已经全面压过同量级的 ResNet，建议把主干换掉。你换了，同样的训练预算、同样的增强，测试准确率掉了七个点。

论文没有说谎，你的实验也没有做错。两个结论都成立，因为它们说的不是同一件事——论文里的数据量是你的一万倍。这一篇要把这个现象讲成一个可以事先估算的权衡，而不是``哪个架构更先进''。

\subsection{卷积把三条先验做成了硬约束}

卷积往架构里塞进三样东西。局部性：一个输出只连到输入的一小片邻域。平移等变：同一片邻域移到别处，用同一条规则处理。参数共享：所有位置共用一组权重。

把它们翻译成对参数空间的操作，力度才看得清。一维长度 \(T\) 的输入接全连接层，权重矩阵有 \(T^2\) 个自由参数，它在 \(\R^{T\times T}\) 里可以任意取值。换成宽度 \(k\) 的卷积，局部性把每行里除了 \(k\) 个之外的元素全部强制为 0，参数共享再把剩下的元素按对角线强制相等——最终矩阵被限制在 \(\R^{T\times T}\) 的一个 \(k\) 维线性子空间里。\(T=512\)、\(k=3\) 时，从 26 万维掉到 3 维。

关键在于``强制''二字。这不是在损失函数上加一项惩罚。Ridge 和 Lasso 那类正则把不喜欢的解压低但不禁止（见\hyperref[sec:13-Machine-Learning/02-Linear-Models/02]{``正则化从 Ridge 到 Lasso''}），数据足够多、证据足够强时，似然项可以压倒惩罚项把解拉出去。架构约束不给这个机会：子空间外的映射根本不在参数化的值域里，梯度下降走不出去，数据再多也走不出去。归纳偏置是先验的极端形式——先验概率在子空间外恒等于零。

\subsection{偏置买到了方差，卖掉的是可表达性}

风险的经典分解把这笔交易摆得很清楚。记 \(\mathcal{H}\) 为假设类，\(f^\star\) 为最优预测函数，\(\hat f_n\) 为在 \(n\) 个样本上训出来的假设。总的超额风险拆成两块：

\[
R(\hat f_n)-R(f^\star)
= \underbrace{\inf_{f\in\mathcal H}R(f)-R(f^\star)}_{\text{近似误差}}
+ \underbrace{R(\hat f_n)-\inf_{f\in\mathcal H}R(f)}_{\text{估计误差}} .
\]

两块对 \(\mathcal{H}\) 的大小反向反应。一致收敛给出的估计误差上界形如 \(\sqrt{\text{复杂度}(\mathcal H)/n}\)，随 \(\mathcal H\) 变小而变小、随 \(n\) 增大而衰减（\hyperref[sec:13-Machine-Learning/10-Learning-Theory/03]{``VC 维把无限假设类变有限''}给出复杂度的一种度量）。近似误差则随 \(\mathcal H\) 变小而变大，且完全不依赖 \(n\)——它是 \(f^\star\) 到类的距离，样本再多也填不平。

条件对账要说清这个分解在深度学习里的适用程度。上界是上界，不是实际风险；深度网络的复杂度度量普遍给出空洞的数值，这一点\hyperref[sec:13-Machine-Learning/10-Learning-Theory/04]{``泛化界的解释力与边界''}讲过。这里用它不是为了算数，而是为了拿到两个定性事实：估计误差随 \(n\) 单调衰减而近似误差不随 \(n\) 变化。仅凭这两条，两条测试误差曲线的形状就定下来了。

\begin{center}
\begin{tikzpicture}[x=1.45cm,y=2.0cm,>=stealth]
\draw[->] (0,0) -- (6.6,0) node[below,font=\scriptsize] {训练数据量（对数尺度）};
\draw[->] (0,0) -- (0,2.15) node[right,font=\scriptsize] {测试误差};
\draw[dotted] (0,0.28) -- (6.4,0.28);
\draw[dotted] (0,0.12) -- (6.4,0.12);
\draw[very thick] plot[domain=0.12:6.3,samples=80] (\x,{0.28+0.95*exp(-\x/1.0)});
\draw[very thick,dashed] plot[domain=0.12:6.3,samples=80] (\x,{0.12+1.75*exp(-\x/2.2)});
\draw[dotted] (5.2,0) -- (5.2,0.29);
\fill (5.2,0.286) circle (1.8pt);
\node[below,font=\scriptsize] at (5.2,-0.03) {交叉点};
\draw[very thick] (3.5,1.72) -- (4.1,1.72);
\node[right,font=\scriptsize] at (4.2,1.72) {强偏置：起点低，早早触顶};
\draw[very thick,dashed] (3.5,1.40) -- (4.1,1.40);
\node[right,font=\scriptsize] at (4.2,1.40) {弱偏置：起点高，持续下降};
\draw[<->] (2.6,0.125) -- (2.6,0.275);
\node[right,font=\scriptsize] at (2.7,0.20) {近似误差之差};
\end{tikzpicture}
\end{center}

\emph{两条曲线各自的水平渐近线是近似误差，与渐近线之间的落差是估计误差；交叉点的位置由这两项的相对大小决定，不由架构的新旧决定。}

图里实线是卷积。假设类小，估计误差衰减得快，所以在小数据量上就已经接近自己的极限；但那条渐近线不低——凡是 \(f^\star\) 里不满足局部性与平移共享的成分，都留在近似误差里，永远消不掉。虚线是弱偏置模型。类大，估计误差起初主导，小数据量上表现难看；随着 \(n\) 增长它一路下压，最终停在一条更低的渐近线上。两条曲线交叉是必然的，唯一的问题是交叉点落在你的数据量的哪一侧。

你的 8000 张缺陷图落在交叉点左边。论文里的数据量落在右边。

\subsection{交叉点可以被挪动}

交叉点不是常数，三件事会改变它的位置。

第一是 \(f^\star\) 偏离先验多远。缺陷检测里，一道划痕出现在工件的哪个部位，识别规则几乎一样，局部纹理就能定性——先验几乎是对的，卷积的近似误差本来就低，交叉点被推得很远。换成需要判断``图左端的物体和右端的物体是不是同一类''的任务，或者绝对位置携带语义的场景（版面分析里页眉与正文的区别、影像里器官的固定解剖位置），平移共享直接与任务冲突，卷积的渐近线被抬高，交叉点大幅左移。

第二是预训练。在别处的大数据上先训一遍，等于用外部样本把估计误差提前压下去，效果是把弱偏置那条曲线整体向左平移。ViT 在大规模预训练后于中等规模下游任务上翻盘，机制就是这个：它自己的下游数据仍在交叉点左边，但曲线已经被搬过去了。这也解释了为什么同一个 ViT 从随机初始化开始训 8000 张图仍然打不过 ResNet——搬曲线的那一步没做。

第三是把偏置以软形式加回去。相对位置编码、局部注意力窗口、卷积词元嵌入，都是在弱偏置模型里重新引入局部性与平移结构，但用可学习的、能被数据推翻的方式，而不是硬约束。这类混合体在中等数据量上通常最稳——它拿到了偏置带来的方差削减，同时保留了数据充足时把偏置推翻的余地。

\subsection{把对比说准}

有两处常见的话术滑坡值得堵住。

一是把 ViT 说成``没有归纳偏置''。它有，只是更弱。切 patch 这一步本身就假设了同一小块内的像素属于同一个语义单元，这是局部性的残留；位置编码给了模型关于空间排列的信息，是位置结构的软形式；多头注意力共享同一组投影矩阵作用于所有位置，那还是参数共享，只不过共享的是``怎么算相似度''而不是``怎么聚合邻域''。真正的差别在于卷积把``和谁交互''钉死在几何邻域上，而注意力让它由内容决定（见\hyperref[sec:14-Deep-Learning/05-Attention-and-Sequence-Models/01]{``Self-Attention 是内容决定的加权读取''}）。这是偏置强弱之别，不是有无之别。

二是把两个架构在基准上的全部差距都记到归纳偏置头上。同期的实验反复表明，把现代训练配方——更长的调度、更重的增强、更好的正则、更合适的归一化——套回卷积主干，相当一部分差距会被抹掉。测得的差异里混着架构与训练协议两个变量，不控制后者就下架构结论，是把配方的功劳算给了先验。

真正能带走的判断是这样一句：选架构等于选先验强度，而先验强度该按``数据量''与``真实结构偏离先验多远''这两个量来定，不按发表时间来定。数据少、任务结构与先验吻合，选强偏置；数据多或者任务里确有长程与非平移结构，选弱偏置并接受它对数据量的胃口；介于两者之间且能拿到预训练权重，选混合体。归纳偏置本身没有好坏，只有匹配与不匹配——这与\hyperref[sec:13-Machine-Learning/03-Trees-Ensembles-and-Neighbors/03]{``KNN 把度量直接变成归纳偏置''}里那个更简单的例子是同一件事，只是这里的偏置被编译进了算子而不是写在距离函数里。

卷积单元到此为止。它的全部力量来自把``和谁交互''固定成几何邻域，全部局限也来自同一处。把这个决定交给内容而不是坐标，是下一个单元的起点。
